\chapter{Acute Pain}

\section*{Definition}
Acute pain is a sensory and emotional experience that usually begins suddenly or recently, is commonly associated with actual or threatened tissue injury or illness, and generally lasts less than 3 months. It may be nociceptive, inflammatory, neuropathic, or mixed, and it often serves as a warning signal. Acute pain can persist or recur and should be reassessed if it does not improve within the expected healing period.

\section*{When to See a Doctor}

\begin{itemize}
 \item Pain is severe, sudden, rapidly worsening, or accompanied by fever, fainting, confusion, weakness, low blood pressure, pale or clammy skin, or other signs of shock
 \item Pain is associated with chest pressure, shortness of breath, sweating, new neurologic symptoms, or a severe sudden headache
 \item Abdominal pain occurs with rigidity, guarding, repeated vomiting, blood in vomit or stool, or inability to pass stool or gas
 \item Pain follows major trauma or occurs with a deformity, open wound, uncontrolled bleeding, loss of function, or inability to bear weight
 \item Severe eye pain with vision changes, sudden testicular pain or swelling, or severe pelvic pain with possible pregnancy requires urgent assessment
\end{itemize}

\section*{How It Develops}
Acute nociceptive pain arises when peripheral nociceptors detect mechanical, thermal, or chemical stimuli and transmit signals through the spinal cord to brain regions involved in location, intensity, and emotional response. Tissue inflammation can release mediators such as prostaglandins and bradykinin that increase sensitivity. Nerve injury can produce neuropathic pain, and the experience of pain is also influenced by sleep, stress, prior pain, expectations, and other biological, psychological, and social factors.

\section*{Causes}
\begin{itemize}
 \item \textbf{Traumatic:} Lacerations, fractures, burns, contusions, crush injuries, sprains, and strains
 \item \textbf{Inflammatory or infectious:} Infection, arthritis, tendinitis, bursitis, abscess, cellulitis, and inflammatory bowel disease
 \item \textbf{Visceral or internal-organ:} Renal colic, biliary colic, myocardial ischemia, pancreatitis, appendicitis, obstruction, peritonitis, and ischemia
 \item \textbf{Post-surgical or procedural:} Incisional pain, tissue retraction, drain sites, and pain after medical or dental procedures
 \item \textbf{Neuropathic:} Acute herpes zoster, nerve entrapment, radiculopathy, and traumatic nerve injury
 \item \textbf{Vascular, neurologic, or metabolic:} Aortic emergencies, stroke-related pain syndromes, migraine, acute glaucoma, sickle-cell pain, and diabetic ketoacidosis
\end{itemize}

\section*{Related Symptoms}
\begin{itemize}
 \item Tenderness, swelling, erythema, muscle spasm, guarding, reduced movement, or loss of function
 \item Fever, chills, sweating, nausea, vomiting, dizziness, weakness, or fainting
 \item Numbness, tingling, weakness, vision changes, chest symptoms, or changes in bowel or bladder function
\end{itemize}
\textbf{Important:} Pain is generally considered chronic when it persists or recurs for longer than 3 months, or beyond the usual healing time. Persistent pain warrants reassessment rather than simply continued self-treatment.

\section*{Medical Evaluation}
\begin{itemize}
 \item Detailed history includes onset, location, quality, severity and effect on function, duration, radiation, aggravating or relieving factors, associated symptoms, injury or procedure, previous episodes, medical history, pregnancy possibility when relevant, medicines, allergies, and alcohol or substance use.
 \item Physical examination is guided by the suspected source and includes vital signs, general appearance, inspection, palpation, range of motion, circulation, and neurologic assessment. Pelvic, rectal, eye, chest, or cardiac examination may be needed.
 \item Targeted laboratory studies may include a complete blood count, electrolytes, kidney and liver tests, inflammatory markers, urinalysis, pregnancy testing, lipase, glucose or ketones, cultures, or troponin according to clinical suspicion.
 \item Imaging is selected by the suspected condition and may include radiography, ultrasound, CT, or MRI. Tests should not be delayed when a time-sensitive emergency is suspected; pregnancy, kidney function, and radiation or contrast risks are considered.
\end{itemize}

\section*{Treatment Options}
\begin{itemize}
 \item Pharmacologic: Non-opioid medicines such as acetaminophen, topical NSAIDs, or oral NSAIDs are often effective, depending on the condition and contraindications. Opioids may be considered for severe acute pain when expected benefits outweigh risks, using immediate-release medicine at the lowest effective dose for no longer than needed.
 \item Non-pharmacologic: Depending on the cause, options include ice, heat, elevation, immobilization, exercise or rehabilitation, relaxation, and other behavioral approaches. Evidence and suitability vary by condition.
 \item Interventional: Nerve blocks, joint injections, regional anesthesia, or other procedures may be appropriate for selected conditions and should be performed by trained clinicians.
\end{itemize}

\section*{Self-Care and Home Management}
\begin{itemize}
 \item For mild pain with no warning signs, use relative rest, ice or heat, elevation, or gentle movement when appropriate for the injury; avoid immobilizing an area longer than advised.
 \item Use over-the-counter acetaminophen or an NSAID only as directed on the label and only if safe for you. Avoid duplicate acetaminophen products and ask a clinician before using NSAIDs if you have kidney disease, ulcers, bleeding risk, take anticoagulants, or are pregnant.
 \item Clean and cover minor wounds, control bleeding with direct pressure, and monitor for increasing redness, swelling, drainage, or fever.
 \item Do not use leftover prescription opioids or someone else\'s medicine. Seek reassessment if pain is severe, worsening, unexplained, or not improving as expected.
\end{itemize}

\section*{References}
\begin{thebibliography}{99}
\bibitem{acute_pain:ref1} International Association for the Study of Pain. Acute Pain; Terminology and Definitions. Current online resources.
\bibitem{acute_pain:ref2} Dowell D, Ragan KR, Jones CM, Baldwin GT, Chou R. CDC Clinical Practice Guideline for Prescribing Opioids for Pain--United States, 2022. \textit{MMWR Recomm Rep}. 2022;71(3):1--95.
\bibitem{acute_pain:ref3} American College of Physicians and American Academy of Family Physicians. Noninvasive Treatments for Acute, Subacute, and Chronic Low Back Pain. \textit{Ann Intern Med}. 2017;166(7):493--505.
\bibitem{acute_pain:ref4} World Health Organization. WHO Guidelines for the Pharmacological and Radiotherapeutic Management of Cancer Pain in Adults and Adolescents. 2019.
\bibitem{acute_pain:ref5} National Library of Medicine. MedlinePlus: Pain and Abdominal Pain. Current online resources.
\end{thebibliography}
